% (User input window)

\begin{theorem}[对角约束最小互信息泛函的 Schur--凹性]\label{thm:pom-diagonal-rate-schur-concavity}
固定有限字母表 \(X=\{1,\dots,n\}\)（\(n\ge 2\)）与 \(\delta\in[0,1]\)。
对任意全支撑分布 \(w\in\Delta(X)\) 定义
\[
R_w(\delta):=\inf\Bigl\{I_P(X;Y):\ P_X=P_Y=w,\ \mathbb P_P(X=Y)\ge 1-\delta\Bigr\}.
\]
则映射 \(w\mapsto R_w(\delta)\) 在 \(\Delta(X)\) 的内点上为 Schur--凹：
对任意 \(v,w\in\Delta(X)\)，若 \(v\prec w\)（Hardy--Littlewood--P\'olya 主序），则
\[
R_v(\delta)\ \ge\ R_w(\delta).
\]
\end{theorem}
\begin{proof}
记 \(p_2(w):=\sum_{i=1}^n w_i^2\)。
若 \(\delta\ge 1-p_2(w)\)，则独立耦合 \(P=w\otimes w\) 满足
\(
\mathbb P(X=Y)=p_2(w)\ge 1-\delta
\)
且 \(I(X;Y)=0\)，从而 \(R_w(\delta)=0\)（命题 \ref{prop:pom-diagonal-rate-zero-threshold-collision}）。
由于 \(t\mapsto t^2\) 凸，\(p_2\) 为 Schur--凸；因此 \(v\prec w\Rightarrow p_2(v)\le p_2(w)\Rightarrow \delta\ge 1-p_2(v)\)，并得到 \(R_v(\delta)=0\)。
故以下仅需讨论 \(\delta<1-p_2(w)\) 的情形。
\par
当 \(\delta<1-p_2(w)\) 时，定理 \ref{thm:pom-diagonal-rate-explicit-param-and-derivative} 给出唯一最优耦合
\[
P^\star_\delta(i,j)=\frac{u_i u_j\bigl(1+\kappa\,\mathbf 1_{\{i=j\}}\bigr)}{Z},
\qquad
w_i=\sum_j P^\star_\delta(i,j)=\frac{u_i(A+\kappa u_i)}{Z},
\qquad
A:=\sum_{j=1}^n u_j,
\]
其中 \(\kappa>0\)、\(u\in(0,\infty)^n\)、\(Z>0\)，且对角质量约束饱和。
由定理 \ref{thm:pom-diagonal-rate-frechet-derivative-scaling-potential}，函数 \(w\mapsto R_w(\delta)\) 在内点处 Fr\'echet 可微；
其切空间梯度可取
\[
\nabla R_w(\delta)\equiv 2\psi_\delta \pmod{\RR\mathbf 1},
\qquad
\psi_\delta(i):=\log\frac{u_i}{w_i}.
\]
\par
由边缘自洽关系 \(\kappa u_i^2+A u_i=Zw_i\) 的正根表示，
\[
\frac{u_i}{w_i}
=\frac{2Z}{A+\sqrt{A^2+4\kappa Zw_i}},
\]
右端随 \(w_i\) 严格递减，因而 \(w_i>w_j\Rightarrow \psi_\delta(i)<\psi_\delta(j)\)。
\par
最后应用 Hardy--Littlewood--P\'olya 的梯度判别准则：对置换不变的可微函数 \(F\)，
\(
F
\)
为 Schur--凹当且仅当对任意 \(i,j\) 有
\[
(w_i-w_j)\Bigl(\frac{\partial F}{\partial w_i}-\frac{\partial F}{\partial w_j}\Bigr)\le 0.
\]
将 \(F(w)=R_w(\delta)\) 代入并用 \(\nabla R_w(\delta)\equiv 2\psi_\delta\pmod{\RR\mathbf 1}\) 消去常数项，即得
\(
(w_i-w_j)(\psi_\delta(i)-\psi_\delta(j))\le 0
\)；
由上一步严格单调性该判别式成立，从而 \(w\mapsto R_w(\delta)\) 为 Schur--凹。证毕。
\end{proof}

\begin{corollary}[均匀分布给出全局上界]\label{cor:pom-diagonal-rate-uniform-global-upper}
令 \(|X|=n\ge 2\)，并记均匀分布 \(U_n\)。
则对任意全支撑分布 \(w\in\Delta(X)\) 与任意 \(\delta\in[0,1]\) 成立
\[
R_w(\delta)\ \le\ R_{U_n}(\delta).
\]
并且当 \(0\le \delta<1-\tfrac1n\) 时，等号当且仅当 \(w\) 为 \(U_n\) 的置换。
此外，\(R_{U_n}(\delta)\) 由定理 \ref{thm:pom-diagonal-rate-uniform-closed-form} 给出显式闭式，等价熵型重写见推论 \ref{cor:pom-diagonal-rate-uniform-entropy-rewrite}。
\end{corollary}
\begin{proof}
均匀分布被任何 \(w\in\Delta(X)\) 主序化，即 \(U_n\prec w\)。
由定理 \ref{thm:pom-diagonal-rate-schur-concavity} 的 Schur--凹性得
\(
R_{U_n}(\delta)\ge R_w(\delta)
\)。
当 \(0\le\delta<1-\tfrac1n\) 时，\(R_{U_n}(\delta)>0\) 且严格性与等号情形由定理 \ref{thm:pom-diagonal-rate-strict-schur-concavity} 给出。证毕。
\end{proof}

\begin{theorem}[严格 Schur--凹性与等号情形]\label{thm:pom-diagonal-rate-strict-schur-concavity}
固定 \(n\ge 2\) 与 \(\delta\in(0,1)\)，并令 \(X=\{1,\dots,n\}\)。
取 \(v,w\in\Delta(X)\) 全支撑且满足
\[
\delta<1-p_2(v),
\qquad
\delta<1-p_2(w),
\]
从而 \(R_v(\delta)>0\)、\(R_w(\delta)>0\)。
若 \(v\prec w\) 且 \(v\) 不是 \(w\) 的置换，则必有严格不等式
\[
R_v(\delta)\ >\ R_w(\delta).
\]
另一方面，若 \(\delta\ge 1-p_2(w)\)，则 \(R_w(\delta)=0\)，并且对任何 \(v\prec w\) 亦有 \(R_v(\delta)=0\)。
\end{theorem}
\begin{proof}
当 \(\delta\ge 1-p_2(w)\) 时，由定理 \ref{thm:pom-diagonal-rate-schur-concavity} 的证明已知 \(R_w(\delta)=0\) 且 \(v\prec w\Rightarrow R_v(\delta)=0\)。
\par
以下设 \(\delta<1-p_2(w)\) 且 \(w\) 全支撑。
由定理 \ref{thm:pom-diagonal-rate-explicit-param-and-derivative} 与定理 \ref{thm:pom-diagonal-rate-frechet-derivative-scaling-potential}，
在内点处
\(
\nabla R_w(\delta)\equiv 2\psi_\delta\pmod{\RR\mathbf 1}
\)
且 \(\psi_\delta(i)=\log(u_i/w_i)\)。
由
\[
\frac{u_i}{w_i}=\frac{2Z}{A+\sqrt{A^2+4\kappa Zw_i}}
\]
可知当 \(w_i\ne w_j\) 时有严格不等式 \(\psi_\delta(i)\ne \psi_\delta(j)\)，并满足同序反向性 \(w_i>w_j\Rightarrow \psi_\delta(i)<\psi_\delta(j)\)。
因此对任意 \(i\ne j\) 且 \(w_i\ne w_j\)，Hardy--Littlewood--P\'olya 的梯度判别式严格成立：
\[
(w_i-w_j)\Bigl(\frac{\partial R_w(\delta)}{\partial w_i}-\frac{\partial R_w(\delta)}{\partial w_j}\Bigr)
\equiv 2(w_i-w_j)\bigl(\psi_\delta(i)-\psi_\delta(j)\bigr)\ <\ 0.
\]
由严格梯度判别准则，\(w\mapsto R_w(\delta)\) 在该区间上为严格 Schur--凹；
于是若 \(v\prec w\) 且非置换，则严格不等式 \(R_v(\delta)>R_w(\delta)\) 成立。证毕。
\end{proof}

\begin{proposition}[最优耦合的二阶子式符号刚性]\label{prop:pom-diagonal-rate-optimal-coupling-minor-sign-rigidity}
设 \(\delta<1-p_2(w)\) 且 \(w\in\Delta(X)\) 全支撑，令 \(P^\star_\delta\) 为实现 \(R_w(\delta)\) 的唯一最优耦合，并写成
\[
P^\star_\delta(i,j)=\frac{u_i u_j\bigl(1+\kappa\,\mathbf 1_{\{i=j\}}\bigr)}{Z},
\qquad
\kappa>0,\ u\in(0,\infty)^n.
\]
则对任意两行指标 \(i\ne j\) 与两列指标 \(k\ne \ell\)，有
\[
\det\!\begin{pmatrix}
P^\star_\delta(i,k) & P^\star_\delta(i,\ell)\\
P^\star_\delta(j,k) & P^\star_\delta(j,\ell)
\end{pmatrix}
=\frac{u_i u_j u_k u_\ell}{Z^2}\Bigl(b_{ik}b_{j\ell}-b_{i\ell}b_{jk}\Bigr),
\qquad
b_{ab}:=1+\kappa\,\mathbf 1_{\{a=b\}}.
\]
特别地，若 \(\{i,j\}\cap\{k,\ell\}=\varnothing\) 则该子式为零；
若交集为单点，则子式的绝对值恒为 \(\kappa\,u_i u_j u_k u_\ell/Z^2\)，其符号由该单一对角元位于主对角角点（\(i=k\) 或 \(j=\ell\)）或反对角角点（\(i=\ell\) 或 \(j=k\)）唯一决定；
若 \(i=k\) 且 \(j=\ell\)，则子式为正且绝对值为 \(\kappa(\kappa+2)\,u_i u_j u_k u_\ell/Z^2\)。
\end{proposition}
\begin{proof}
将 \(P^\star_\delta(i,j)=\frac{u_i u_j}{Z}b_{ij}\) 代入并将 \(u_i,u_j,u_k,u_\ell\) 提出到行列式外即得显式分解。
由于 \(b_{ab}\in\{1,1+\kappa\}\) 仅依赖于指标相等关系，后三条结论为有限情形的直接枚举。证毕。
\end{proof}

